% ROUND6-REFEREE-POSITIVE-CLOSURE
\section{Round-six positive closure: Ray histories, transmission-zero memory, and conditional rough limits}
\label{sec:r6-a4}

The exact state is the complete resolved history.  The finite-dimensional
memory theorem is stated without a false orthogonality property.  Complex
zeros of the compressed correlation transform are retained as finite auxiliary
modes; exponential decay is asserted only for the resulting minimum-phase
residual kernel.  The enhanced invariance principle includes the deterministic
antisymmetric area drift and is strengthened to conditional kernels on compact
history sets.

\subsection{A common-exceptional-set history process}

Let \(E\) be the compact graph completion of the Sinai suspension state and
let
\[
\mathsf H=D(( -\infty,0],E)
\]
with the local Skorokhod sigma-field.  Fix a prepared stationary path law
\(\nu\).  For rational \(t\ge0\), disintegrate the shifted future given the
past and denote the kernel by \(P_t(h,\cdot)\).

\begin{lemma}[Simultaneous conditional kernels]
\label{lem:r6-a4-kernels}
There is one invariant \(\nu\)-null set \(N\subset\mathsf H\) such that, on
\(\mathsf H\setminus N\), all rational-time kernels are defined,
right-continuous on the cylinder algebra, and satisfy
\[
P_{s+t}=P_sP_t
\]
for all nonnegative rational \(s,t\).  Sending \(N\) to a cemetery history
produces a Markov family on the complete measurable history space.
\end{lemma}

\begin{proof}
Choose a countable convergence-determining cylinder algebra and a countable
set of rational time pairs.  For each conditional expectation and each
Chapman--Kolmogorov identity, select a version and remove its exceptional set.
The countable union is still null.  Stationarity makes the union invariant
after adding all rational shifts.  Right continuity of the graph-completed
path gives right continuity of every cylinder conditional expectation outside
one further countable union of null sets.  Define the cemetery kernel on the
remaining exceptional set.
\end{proof}

Let \(R_\lambda F=\int_0^\infty e^{-\lambda t}P_tFdt\) on the countable
cylinder algebra.  The Ray cone is the rational lattice generated by
\(R_\lambda F\), constants, minima, and positive rational combinations.  Its
evaluation map gives a separable Ray compactification \(\mathsf H_R\).

\begin{theorem}[Measured Ray realization]
\label{thm:r6-a4-ray}
The kernels of Lemma~\ref{lem:r6-a4-kernels} extend to a Borel right process
on the Polish Ray space \(\mathsf H_R\).  The Ray resolvent separates histories
modulo the completed \(\nu\)-null relation, the completed measurable sets
agree with those of the original history process, and the transition
semigroup is strongly continuous in the bounded-strict topology generated by
compact Ray sets.
\end{theorem}

\begin{proof}
The resolvent identity
\[
R_\lambda-R_\mu=(\mu-\lambda)R_\lambda R_\mu
\]
follows from the rational-time semigroup and monotone approximation of the
integral.  The Ray cone is countable and separates the completed history
sigma-field because the original cylinder algebra does.  Its evaluation
closure is therefore compact metrizable on each bounded Ray level and Polish
in the increasing union.  The kernels preserve excessive functions and hence
extend uniquely to the Ray closure.  Resolvent right continuity and the common
exceptional-set construction give a right-continuous strong Markov process.
For a bounded-strict function, approximate uniformly on a compact Ray set by
finite Ray-cone combinations; right continuity is uniform on that compact set,
while the complement is suppressed by the strict multiplier.  This proves
strong continuity.
\end{proof}

On an exposed spectral phase, use the weighted history metric
\[
d_\beta(h,\tilde h)=
\sum_{m\ge1}2^{-m}
\bigl(1\wedge d_{J_1}(h|_{[-m,0]},\tilde h|_{[-m,0]})\bigr).
\]
The A2 spectral gap and the A3 singularity shield couple two histories after
their next common Young return.

\begin{theorem}[Concrete weighted-history Feller theorem]
\label{thm:r6-a4-feller}
For every exposed phase, the Ray realization is equivalent modulo null sets to
the weighted Skorokhod history space.  Its semigroup is a strongly continuous
contraction on \(\mathrm{BUC}_\beta(\mathsf H)\), and for a cylinder depending
on the last \(m\) units,
\[
\operatorname{Lip}_\beta(P_tF)
\le Ce^{-\gamma(t-m)_+}\operatorname{Lip}_\beta(F)+C\|F\|_\infty.
\]
Dynamically Hölder cylinders form a core.
\end{theorem}

\begin{proof}
Shield the two histories from singular neighborhoods and long return blocks.
On the shielded event, the finite connector property gives a common return
word and the quotient spectral gap couples the futures with probability
\(1-Ce^{-\gamma(t-m)_+}\).  A3 makes the unshielded event uniformly small.
This gives the Lipschitz estimate and Feller property.  Time averages
\(n\int_0^{1/n}P_tFdt\) approximate every bounded uniformly continuous
function and lie in the graph closure of the cylinder derivation, proving the
core statement.  The cylinder evaluations identify the Ray and weighted
history realizations modulo the common exceptional set.
\end{proof}

\subsection{The physical history eigenfunction and Doob transform}

Let \(\Psi\) be a bounded dynamically Hölder path potential on an A2 simple
spectral chart.  Denote the leading suspension eigenvalue by \(\lambda_\Psi\),
the positive quotient eigenfunction by \(h_Y\), and the roof by \(\tau\).
For a history \(h\), let \(y(h)\) be its last Young base state and let
\(s(h)\) be the elapsed time since that return.  Define
\[
h_\Psi(h)=
 e^{-\lambda_\Psi s(h)}
\mathbb E_{y(h)}\left[
 e^{\int_0^{s(h)}\Psi(H_u)du}
 h_Y(Y_{s(h)})
\right].
\]

\begin{theorem}[History eigen-triple]
\label{thm:r6-a4-eigen}
The definition above is independent of coding representatives at
one-sided graph-completed singular histories, extends to a strictly positive
bounded function with bounded inverse on the weighted history space, and
satisfies
\[
P_t^\Psi h_\Psi=e^{\lambda_\Psi t}h_\Psi.
\]
Consequently
\[
\widetilde P_t^\Psi F
=e^{-\lambda_\Psi t}h_\Psi^{-1}P_t^\Psi(h_\Psi F)
\]
is a strongly continuous Markov semigroup on
\(\mathrm{BUC}_\beta(\mathsf H)\).
\end{theorem}

\begin{proof}
The quotient eigenfunction has matching one-sided traces in the A2 ambient
trace field.  Thus the two coding representatives at a singular history have
the same graph-completed value.  Bounded roof and bounded potential give
upper and lower bounds during one return; the quotient eigenfunction is also
bounded above and away from zero by the spectral cone theorem.  The Markov
property at the next return, followed by the quotient eigenrelation, proves
the identity first at return times and then at every physical time by the
suspension formula.  Multiplication by \(h_\Psi^{\pm1}\) preserves the
weighted history space, so strong continuity follows from
Theorem~\ref{thm:r6-a4-feller}.
\end{proof}

For \(\theta\ne0\), the nonlinear history pressure semigroup is
\[
\mathcal E_t^{\Psi,\theta}F
=\theta^{-1}\log\widetilde P_t^\Psi(e^{\theta F}).
\]
It has the exact nonlinear tower by the Markov property of
\(\widetilde P_t^\Psi\).

\subsection{Domain-safe compressed memory}

Work in \(\mathcal H=L^2(\nu_\Psi^D)\), where \(\nu_\Psi^D\) is the invariant
Doob phase.  Let \(U_t\) be the Koopman or Markov contraction semigroup with
generator \(L\), let \(V\subset D(L^2)\) be finite dimensional, and let
\(P\) be the orthogonal projection onto \(V\).  Put
\[
C(t)=PU_tP|_V,
\qquad A=C'(0)=PLP|_V.
\]
Since \(C(0)=I_V\), there is a unique continuous kernel \(K\) solving
\[
C'(t)=AC(t)+(K*C)(t).
\]
For \(B\in D(L)\), set \(x_B(t)=PU_tB\) and
\[
\eta_B(t)=x_B'(t)-Ax_B(t)-(K*x_B)(t).
\]

\begin{theorem}[Exact Volterra identity and correct forcing trace]
\label{thm:r6-a4-volterra}
The kernel \(K\) is locally continuous and has Laplace transform
\[
\widehat K(z)=zI-A-[P(z-L)^{-1}P|_V]^{-1},
\qquad \Re z>0.
\]
For every \(B\in D(L)\),
\[
x_B'(t)=Ax_B(t)+(K*x_B)(t)+\eta_B(t)
\]
holds exactly.  At time zero,
\[
\eta_B(0)=PLQB,
\]
which is a resolved vector and is not asserted to be orthogonal to \(V\).
If \(B\in V\), then \(\eta_B\equiv0\).
\end{theorem}

\begin{proof}
Differentiate \(C(t)\) twice using \(V\subset D(L^2)\).  The second-kind
Volterra equation has a unique continuous solution by Picard iteration on
compact intervals.  Taking Laplace transforms gives the displayed formula.
The residual definition yields the exact equation for every \(B\).  At zero,
\(x_B'(0)=PLB\), \(x_B(0)=PB\), and the convolution vanishes, hence
\(\eta_B(0)=PLB-PLPB=PLQB\).  For \(B\in V\), uniqueness of the Volterra
solution for \(C(t)B\) gives zero residual.
\end{proof}

\subsection{Transmission zeros and an augmented minimum-phase realization}

The A2 spectral packet makes every entry of
\[
\widehat C(z)=P(z-L)^{-1}P|_V
\]
meromorphic in a left half-strip and gives
\(\widehat C(z)=z^{-1}I+O(z^{-2})\) at large vertical frequency.  Hence every
closed smaller half-strip contains finitely many poles and finitely many zeros
of \(\det\widehat C\), counted with multiplicity.

\begin{lemma}[Finite transmission-zero realization]
\label{lem:r6-a4-zeros}
Fix \(0<\gamma_1<\gamma_0\).  There is a finite-dimensional auxiliary space
\(Z\), matrices \((A_Z,B_Z,C_Z,D_Z)\), and an enlarged resolved state
\(\widetilde V=V\oplus Z\) such that:
\begin{enumerate}[label=(\roman*)]
\item the input-output transfer on the original \(V\)-coordinate is exactly
\(\widehat C(z)\);
\item every zero of \(\det\widehat C\) in
\(\Re z\ge-\gamma_1\) appears as an eigenmode of \(A_Z\);
\item the compressed transfer of the enlarged state is analytic and
invertible on \(\Re z\ge-\gamma_1\), apart from the spectral poles already
placed in the finite instantaneous matrix.
\end{enumerate}
\end{lemma}

\begin{proof}
Take the Smith--McMillan factorization of the finite matrix meromorphic
function on the strip.  Its finitely many zero divisors form a proper rational
matrix factor \(B(z)\); realize \(B(z)^{-1}\) by a finite controllable and
observable state-space system on \(Z\).  Multiply the original transfer by
this inverse factor internally and expose the auxiliary state as part of the
resolved coordinate.  Pole factors are treated in the same way.  The
remaining matrix factor has neither poles nor zeros on the closed smaller
strip and tends to \(z^{-1}I\) at infinity, so it and its inverse are analytic
there.  Standard block elimination shows that eliminating \(Z\) recovers the
original transfer exactly.
\end{proof}

\begin{theorem}[Finite modes plus exponentially decaying memory]
\label{thm:r6-a4-decay}
For the augmented state \(\widetilde V\), the residual memory kernel
\(K_{\rm res}\) satisfies
\[
\|K_{\rm res}(t)\|\le Ce^{-\gamma t}
\]
for some \(\gamma>0\).  The original projected dynamics is therefore exactly
represented by a finite system containing all resonance and transmission-zero
modes, coupled to an exponentially integrable residual memory.  No zero-free
claim is deduced from covariance positivity.
\end{theorem}

\begin{proof}
By Lemma~\ref{lem:r6-a4-zeros}, the augmented memory transform is analytic in
a left half-strip.  Repeated resolvent differentiation and the A2
high-frequency estimate give an integrable vertical majorant after finitely
many integrations by parts.  Shift the Bromwich contour to
\(\Re z=-\gamma\).  The finite residues have already been placed in the
instantaneous auxiliary matrix, leaving the stated exponential bound.  Block
elimination recovers the original projected coordinate.
\end{proof}

\subsection{Enhanced and conditional invariance principles}

Let \(G\) be a centered vector of dynamically Hölder suspension observables
with a moment exponent \(q>4\).  Put
\[
X_t^\varepsilon=\varepsilon
\int_0^{t/\varepsilon^2}G\circ\Phi^sds.
\]
Define
\[
D=\int_{-\infty}^{\infty}
\operatorname{Cov}(G,G\circ\Phi^t)dt,
\]
\[
\Gamma=\frac12\int_0^\infty
\left[
\operatorname{Cov}(G,G\circ\Phi^t)
-\operatorname{Cov}(G\circ\Phi^t,G)
\right]dt.
\]

\begin{theorem}[Area-corrected enhanced invariance principle]
\label{thm:r6-a4-rough}
For every \(p>2\), the canonical lift
\((X^\varepsilon,\mathbb X^\varepsilon)\) converges in
\(p\)-variation to
\[
\left(B,
\int(B-B_s)\otimes\circ dB+(t-s)\Gamma
\right),
\]
where \(B\) has covariance \(D\).  The normalization and sign of
\(\Gamma\) are those displayed above.
\end{theorem}

\begin{proof}
The A2 spectral gap gives a vector martingale--coboundary decomposition with
an \(L^q\) martingale difference and an exponentially summable transfer
remainder.  The martingale rough invariance principle gives the Brownian first
level and its Stratonovich second level.  Summing the cross terms between the
martingale and coboundary parts produces the antisymmetric correlation sum;
roof interpolation converts the discrete sum to the displayed integral.
Burkholder bounds of order \(q>4\) give tightness in every
\(p\)-variation topology with \(p>2\).  The symmetric correction is absorbed
in the Stratonovich lift, leaving \(\Gamma\).
\end{proof}

\begin{theorem}[Uniform conditional-kernel homogenization]
\label{thm:r6-a4-conditional}
Let \(K_T^\varepsilon(h,\cdot)\) be the conditional law of the enhanced slow
path over a fixed macroscopic horizon, given the complete history \(h\).
For every compact weighted-history set \(K\) and every bounded Lipschitz
functional of a finite enhanced path window,
\[
\sup_{h\in K}
\left|K_T^\varepsilon F(h)-K_T^{\rm diff}F(X_0(h))\right|
\longrightarrow0,
\]
where the limiting rough diffusion includes the area drift \(\Gamma\).  The
same convergence holds stably with any bounded variable measurable with
respect to the initial history.
\end{theorem}

\begin{proof}
Truncate the remote past at length \(m\).  The coupling estimate in
Theorem~\ref{thm:r6-a4-feller} makes the conditional-kernel error
\(Ce^{-\gamma m}\), uniformly on \(K\).  For a fixed truncated history, the
A2 transfer operator with a central finite-block insertion gives uniform
characteristic-function convergence of the martingale array and all second
iterated sums.  The \(L^q\) bounds used in
Theorem~\ref{thm:r6-a4-rough} are uniform over the compact truncated-history
family, so the convergence is uniform and stable.  First let
\(\varepsilon\to0\), then \(m\to\infty\).
\end{proof}

\subsection{The coupled slow-memory limit}

Let \(Y^\varepsilon\) solve the rough differential equation
\[
dY_t^\varepsilon=V_0(Y_t^\varepsilon)dt+
V(Y_t^\varepsilon)dX_t^\varepsilon+
\int_0^tK_{\rm res}^\varepsilon(t-s)R(Y_s^\varepsilon)ds\,dt
+C_ZZ_t^\varepsilon dt,
\]
with the auxiliary transmission-zero state
\[
dZ_t^\varepsilon=A_ZZ_t^\varepsilon dt+B_ZR(Y_t^\varepsilon)dt.
\]
Assume compatible stationary initial unresolved histories and
\(K_{\rm res}^\varepsilon(t)=\varepsilon^{-2}
K_{\rm res}(t/\varepsilon^2)\).

\begin{theorem}[Round-six A4 closure]
\label{thm:r6-a4-main}
Prepared Sinai paths possess a common-exceptional-set right history process,
a concrete exposed-phase Feller realization, a genuine eigenfunction Doob
semigroup, an exact domain-safe Volterra equation, and a finite-mode plus
exponentially decaying memory realization which retains every transmission
zero.  The enhanced slow limit contains the deterministic area anomaly, and
the conditional kernels converge uniformly on compact history sets.  The
coupled slow-memory system converges by continuity of the rough Volterra map
to the corresponding diffusion with the area-induced drift and the finite
auxiliary modes.
\end{theorem}

\begin{proof}
The history assertions follow from Theorems~\ref{thm:r6-a4-ray},
\ref{thm:r6-a4-feller}, and \ref{thm:r6-a4-eigen}.  The memory assertions
follow from Theorems~\ref{thm:r6-a4-volterra} and
\ref{thm:r6-a4-decay}.  The enhanced and conditional limits are
Theorems~\ref{thm:r6-a4-rough} and
\ref{thm:r6-a4-conditional}.  Exponential integrability of the residual
kernel gives a uniform first moment, so the rescaled convolution converges to
its mass on tight rough paths.  The rough Volterra solution map is locally
Lipschitz on bounded sets; joint convergence with the finite linear auxiliary
state completes the proof.
\end{proof}
